\chapter{Lezione 4 --- Teorema del campionamento, banda pratica, rumore, filtro anti-alias e modello di sampling in Laplace}

\section{Richiamo: campionamento nel tempo e periodicizzazione in frequenza}

Ripartiamo dal modello già introdotto nella lezione precedente.
Un segnale continuo \(f(t)\) viene campionato moltiplicandolo per un pettine di Dirac
di periodo \(T\):
\[
\operatorname{III}_T(t)=\sum_{k=-\infty}^{+\infty}\delta(t-kT).
\]

Il segnale campionato è quindi
\[
\hat f(t)=f(t)\,\operatorname{III}_T(t)
=
f(t)\sum_{k=-\infty}^{+\infty}\delta(t-kT)
=
\sum_{k=-\infty}^{+\infty} f(kT)\,\delta(t-kT).
\]

Dal punto di vista del tempo, questa operazione sostituisce il segnale continuo
con un treno di impulsi pesati con i valori dei campioni.

Nel dominio della frequenza, la trasformata di Fourier del segnale campionato
risulta periodica e ripete lo spettro originale con periodo
\[
\omega_s=\frac{2\pi}{T},
\]
se si lavora nel dominio delle pulsazioni, oppure con periodo
\[
f_s=\frac{1}{T},
\]
se si lavora nella frequenza ordinaria in Hertz.

Più precisamente, se \(F(\omega)\) è la trasformata di Fourier del segnale continuo,
allora
\[
\hat F(\omega)=\frac{1}{T}\sum_{n=-\infty}^{+\infty} F(\omega-n\omega_s).
\]

Quindi il campionamento:
\begin{itemize}
    \item discretizza il segnale nel tempo;
    \item replica periodicamente il suo spettro in frequenza.
\end{itemize}

\section{Ricostruzione ideale del segnale continuo}

Poiché tra dominio del tempo e dominio della frequenza esiste una relazione biunivoca
tramite trasformata e antitrasformata di Fourier, per ricostruire idealmente \(f(t)\)
a partire dal segnale campionato si procede così:
\begin{enumerate}
    \item si considera lo spettro periodico \(\hat F(\omega)\) del segnale campionato;
    \item si isola una sola copia dello spettro originale;
    \item si applica l'antitrasformata di Fourier di quella sola porzione spettrale.
\end{enumerate}

Questo procedimento è possibile solo se le repliche spettrali non si sovrappongono.

\section{Teorema del campionamento}

Supponiamo che il segnale sia a banda limitata, cioè che esista una pulsazione \(\bar\omega\)
tale che
\[
F(\omega)=0
\qquad
\text{per } |\omega|>\bar\omega.
\]

Allora il segnale è ricostruibile dai suoi campioni se la frequenza di campionamento
è sufficientemente alta da impedire la sovrapposizione delle copie periodiche dello spettro.

La condizione fondamentale è
\[
\omega_s > 2\bar\omega.
\]

Poiché
\[
\omega_s=\frac{2\pi}{T},
\]
questa condizione è equivalente a
\[
\frac{2\pi}{T}>2\bar\omega
\qquad \Longleftrightarrow \qquad
T<\frac{\pi}{\bar\omega}.
\]

Nel dominio delle frequenze ordinarie, la stessa condizione si scrive
\[
f_s>\,2\bar f.
\]

\subsection{Interpretazione}

La distanza tra le repliche dello spettro è fissata da \(1/T\) oppure da \(\omega_s\):
\begin{itemize}
    \item se \(T\) è piccolo, le repliche in frequenza sono lontane;
    \item se \(T\) è grande, le repliche in frequenza si avvicinano.
\end{itemize}

Per campionare segnali con frequenze elevate bisogna quindi ridurre il periodo di campionamento,
cioè aumentare la frequenza di campionamento.

\section{Dualismo tempo--frequenza}

Negli appunti viene richiamato un principio qualitativo molto importante:
c'è un dualismo tra localizzazione temporale e localizzazione in frequenza.

Nel caso del campionamento:
\begin{itemize}
    \item campioni molto distanziati nel tempo (cioè \(T\) grande) producono repliche
    molto vicine in frequenza;
    \item campioni molto fitti nel tempo (cioè \(T\) piccolo) producono repliche
    molto lontane in frequenza.
\end{itemize}

Questo spiega geometricamente perché la ricostruzione peggiora quando si campiona troppo lentamente.

\section{Segnali a banda limitata: caso ideale e caso pratico}

Dal punto di vista matematico, un segnale \emph{strettamente} a banda limitata
ha trasformata di Fourier nulla oltre una certa frequenza di taglio.

Tuttavia un segnale davvero a banda limitata è un modello ideale:
in molti casi reali il segnale non si esaurisce perfettamente in frequenza,
ma presenta code spettrali che decadono progressivamente.

Per questo, in pratica, si introduce una \textbf{banda passante effettiva}
o \textbf{banda pratica}, definita tramite una soglia convenzionale.

\section{Definizione pratica di banda: il criterio dei \(-3\) dB}

Negli appunti compare il richiamo al valore \(-3\) dB.

Se \(F(\omega)\) ha un massimo a bassa frequenza, si può definire la banda passante
come l'intervallo di frequenze per cui il modulo resta sopra una soglia prefissata.
Una scelta molto comune è la soglia a \(-3\) dB rispetto al massimo.

In questo modo si definisce una pulsazione di banda \(\bar\omega\) come il punto
in cui il modulo è sceso di \(3\) dB rispetto al valore di riferimento.

\subsection{Osservazione importante}

La soglia dei \(-3\) dB non è una legge assoluta:
\begin{itemize}
    \item è una convenzione pratica molto diffusa;
    \item in altri casi si può usare una soglia diversa, ad esempio \(-6\) dB;
    \item la scelta va fatta in base al sistema e alla precisione desiderata.
\end{itemize}

Quindi, per i segnali reali, la nozione di banda è spesso una nozione \emph{operativa}
più che rigorosamente matematica.

\section{Rumore di misura}

Quando si misura un segnale fisico con un sensore, la misura non è mai esatta.
Si può modellare il segnale misurato come
\[
f_m(t)=f(t)+m(t),
\]
dove:
\begin{itemize}
    \item \(f(t)\) è il segnale ``vero'';
    \item \(m(t)\) è il rumore di misura.
\end{itemize}

Il rumore dipende da:
\begin{itemize}
    \item sensibilità del sensore;
    \item elettronica di misura;
    \item ambiente;
    \item processo di acquisizione.
\end{itemize}

Per questo, nel problema del campionamento, non si ha quasi mai a che fare
con un segnale puro, ma con un segnale contaminato dal rumore.

\section{Rumore bianco}

Negli appunti compare il richiamo al \textbf{processo bianco}.

\subsection{Autocorrelazione}

Un processo di rumore bianco ideale ha autocorrelazione del tipo
\[
R_m(\tau)=\mathbb{E}\{m(t)m(t+\tau)\}=K\,\delta(\tau),
\]
dove \(K\) è una costante positiva.

Questo significa che campioni presi a istanti diversi sono incorrelati.

\subsection{Spettro del rumore bianco}

La trasformata di Fourier della delta è costante, quindi la densità spettrale
del rumore bianco è piatta:
\[
S_m(\omega)=\mathcal{F}\{R_m(\tau)\}=K.
\]

Per questo si dice che un rumore bianco ha energia distribuita su tutte le frequenze.

\subsection{Rumore bianco e rumore gaussiano}

Le due proprietà non coincidono:
\begin{itemize}
    \item \textbf{bianco} significa spettro piatto / campioni incorrelati;
    \item \textbf{gaussiano} si riferisce invece alla distribuzione di probabilità.
\end{itemize}

Un rumore può essere gaussiano ma non bianco, oppure bianco ma non gaussiano.

\section{Perché il rumore crea problemi nel campionamento}

Se il segnale misurato è
\[
f_m(t)=f(t)+m(t),
\]
allora il suo spettro contiene:
\begin{itemize}
    \item la parte utile \(F(\omega)\);
    \item la parte di rumore \(M(\omega)\).
\end{itemize}

Quando si campiona, anche lo spettro del rumore viene replicato periodicamente.
Poiché il rumore bianco è presente su tutte le frequenze, le repliche possono ripiegarsi
sulle basse frequenze e contaminare la banda utile.

Questo fenomeno è una delle ragioni principali per cui, \emph{prima del campionamento},
si introduce un filtro anti-alias.

\section{Filtro anti-alias}

Prima di campionare un segnale reale è quindi necessario inserire un
\textbf{filtro anti-alias}, cioè un filtro analogico posto a monte del campionatore.

Lo schema concettuale è:
\[
f(t)\ \longrightarrow\ \text{filtro anti-alias}\ \longrightarrow\ \text{campionatore}\ \longrightarrow\ \hat f(kT).
\]

\subsection{Funzione del filtro anti-alias}

Il filtro anti-alias deve:
\begin{itemize}
    \item attenuare le componenti ad alta frequenza;
    \item limitare il contenuto spettrale che, dopo il campionamento,
    si ripiegherebbe nella banda utile;
    \item rendere applicabile in modo pratico il teorema del campionamento.
\end{itemize}

Idealmente sarebbe un passa-basso rettangolare, ma nella realtà si usano filtri
con pendenza finita.

\section{Caratteristiche del filtro anti-alias}

Il filtro anti-alias è tipicamente un filtro passa-basso analogico.

Un esempio semplice è un filtro RC del primo ordine, con pendenza asintotica
\[
-20\ \text{dB/decade}.
\]

Filtri di ordine più elevato permettono pendenze maggiori, ma introducono
in generale maggiore complessità e maggiore sfasamento.

\subsection{Compromesso attenuazione--sfasamento}

Un punto essenziale della lezione è che il filtro anti-alias non può essere scelto
guardando solo l'attenuazione.

Infatti:
\begin{itemize}
    \item una pendenza più rapida attenua meglio le alte frequenze;
    \item però una pendenza più rapida introduce più sfasamento;
    \item lo sfasamento può degradare la dinamica del sistema e, in alcuni casi,
    compromettere anche i margini di stabilità.
\end{itemize}

La scelta del filtro è quindi sempre un compromesso tra:
\begin{itemize}
    \item quanto rumore / aliasing si vuole sopprimere;
    \item quanto ritardo di fase si può tollerare.
\end{itemize}

\section{Scelta pratica di frequenza di taglio e frequenza di campionamento}

Nel caso ideale il teorema di campionamento richiederebbe soltanto
\[
\omega_s>2\bar\omega.
\]

Nella pratica però, poiché il filtro anti-alias non è ideale, si preferisce lavorare
con un margine di sicurezza:
\begin{itemize}
    \item si sceglie la frequenza di campionamento ben più alta del doppio della banda utile;
    \item si colloca la frequenza di taglio del filtro anti-alias in modo coerente
    con la banda del segnale e con la pendenza disponibile.
\end{itemize}

Negli appunti è richiamata una regola pratica secondo cui conviene scegliere
\[
\omega_c \gg 2\bar\omega,
\]
oppure, in forma equivalente,
\[
f_s \text{ sufficientemente maggiore di } 2\bar f.
\]

Una regola empirica spesso usata in pratica è quella di scegliere la frequenza di campionamento
diverse volte più grande della banda utile, ad esempio
\[
f_s \ge 5\cdot (2\bar f)
\quad \text{oppure} \quad
f_s \ge 10\cdot (2\bar f),
\]
a seconda delle prestazioni richieste e della pendenza del filtro disponibile.

\subsection{Osservazione}

Se si sbaglia la scelta della frequenza di campionamento oppure della frequenza di taglio
del filtro anti-alias, l'informazione persa per aliasing non può più essere recuperata a valle.

\section{Sensori digitali e sensori analogici}

Molti sensori reali sono ancora analogici, e in tal caso:
\begin{itemize}
    \item il filtraggio anti-alias;
    \item il campionamento;
\end{itemize}
vengono effettuati nel sistema di acquisizione.

Se invece il sensore è digitale, spesso:
\begin{itemize}
    \item il costruttore incorpora già una fase di filtraggio;
    \item il periodo di campionamento è prefissato oppure regolabile in modo limitato.
\end{itemize}

In ogni caso, dal punto di vista progettuale, rimane centrale capire quali siano
le caratteristiche in frequenza del filtro anti-alias effettivo.

\section{Passaggio dal modello in Fourier al modello in Laplace}

Finora il campionamento è stato descritto soprattutto nel dominio della Fourier,
cioè guardando la periodicizzazione dello spettro.

Tuttavia, per l'analisi e il progetto dei sistemi di controllo, è più comodo lavorare
nel dominio di Laplace, perché:
\begin{itemize}
    \item i modelli dei sistemi dinamici continui sono dati in \(s\);
    \item il regolatore e l'impianto vengono normalmente descritti tramite
    funzioni di trasferimento;
    \item il campionamento interessa solo una parte dell'anello di controllo.
\end{itemize}

Per questo si introduce un modello di sampling compatibile con la trasformata di Laplace.

\section{Modello di sampling nel dominio di Laplace}

Per un segnale causale campionato agli istanti \(kT\), si considera
\[
\hat f(t)=\sum_{k=0}^{+\infty} f(kT)\,\delta(t-kT).
\]

Calcoliamo la sua trasformata di Laplace:
\[
\mathcal{L}\{\hat f(t)\}
=
\int_{0^-}^{+\infty}
\left(
\sum_{k=0}^{+\infty} f(kT)\,\delta(t-kT)
\right)e^{-st}\,dt.
\]

Portando fuori la sommatoria,
\[
\mathcal{L}\{\hat f(t)\}
=
\sum_{k=0}^{+\infty} f(kT)\,\mathcal{L}\{\delta(t-kT)\}.
\]

Poiché
\[
\mathcal{L}\{\delta(t-kT)\}=e^{-skT},
\]
si ottiene
\[
\mathcal{L}\{\hat f(t)\}
=
\sum_{k=0}^{+\infty} f(kT)e^{-skT}.
\]

Questa non è una funzione razionale in \(s\), perché contiene l'esponenziale \(e^{-sT}\).

\section{Introduzione della trasformata \(z\)}

Per trattare comodamente il campionamento si introduce la variabile
\[
z=e^{sT}.
\]

Di conseguenza
\[
e^{-sT}=z^{-1}.
\]

Sostituendo nella formula precedente,
\[
\mathcal{L}\{\hat f(t)\}
=
\sum_{k=0}^{+\infty} f(kT)\,z^{-k}.
\]

Questa espressione suggerisce la definizione della \textbf{trasformata \(z\)} del segnale campionato:
\[
\mathcal{Z}\{f[k]\}
=
\sum_{k=0}^{+\infty} f(kT)\,z^{-k}.
\]

Quindi la trasformata \(z\) nasce come strumento naturale per descrivere
segnali e sistemi campionati.

\section{Significato della trasformata \(z\)}

La variabile \(z\) incorpora in sé l'informazione del campionamento,
perché è legata al periodo \(T\) attraverso la relazione
\[
z=e^{sT}.
\]

La trasformata \(z\):
\begin{itemize}
    \item svolge per i sistemi discreti un ruolo analogo a quello della Laplace
    per i sistemi continui;
    \item permette di passare da equazioni alle differenze a funzioni di trasferimento discrete;
    \item è lo strumento naturale per lo studio dei regolatori digitali.
\end{itemize}

\section{Sintesi della lezione}

I punti fondamentali della lezione sono:
\begin{enumerate}
    \item richiamo del teorema del campionamento e della condizione
    \[
    \omega_s>2\bar\omega
    \qquad\text{oppure}\qquad
    T<\frac{\pi}{\bar\omega};
    \]
    \item interpretazione del dualismo tempo--frequenza;
    \item distinzione tra banda ideale e banda pratica;
    \item uso del criterio a \(-3\) dB per definire una banda passante operativa;
    \item introduzione del rumore di misura e del rumore bianco;
    \item necessità del filtro anti-alias prima del campionamento;
    \item compromesso tra attenuazione delle alte frequenze e sfasamento introdotto dal filtro;
    \item passaggio dal modello di sampling in Fourier al modello di sampling in Laplace;
    \item nascita della trasformata \(z\) dalla relazione
    \[
    z=e^{sT}.
    \]
\end{enumerate}